\documentclass[a4paper,12pt]{letter}

\usepackage{amsmath}
\usepackage{amssymb}
\usepackage{nopageno}
\usepackage{amsmath}
\usepackage{enumitem}
% \usepackage{mathtools}
\usepackage{eqnarray,amsmath}
\usepackage[a4paper, total={7.5in, 11.25in}]{geometry}
%\usepackage[margin=2.5cm]{geometry}
\usepackage[utf8]{inputenc}
\usepackage[T1]{fontenc}
\usepackage{lmodern}
%\usepackage{amsmath, amssymb}
\usepackage{setspace}
\setstretch{1.2}

\begin{document}

\begin{center}
 (KTXFI2EBNF) Physics II.  Exercises \\ \smallskip
November~10,~2025
 \end{center}

\begin{enumerate}[label=\Roman*.]

\item Relativistic Kinematic
 
\begin{enumerate}[label=\arabic*.] 

\item A spaceship with a rest length of 100\,m passes in front of an Earth observer in $4\,\mu s$.  What is its velocity relative to the Earth?

\item We study an electron in a radioactive particle beam as it passes through the laboratory.  It is found that, as measured in the laboratory, the average lifetime of the particle is 20 ns; after this time the particle decays.  The average lifetime measured in the reference frame attached to the particles is $7.5\,ns$.  What is the propagation velocity of the particle beam?

\item A certain bacterium doubles its population in 20\,days.  Two such bacteria are placed in a spaceship and sent on a 1000-day space journey.  During the trip, the spaceship moves at a constant speed of $0.995 \, c$. How many bacteria will be on board at the moment of return?

\item Aladár drives a nail into point $ P_1(4000\,\text{m}, -100\,\text{m}, 200\,\text{m}) $ in the coordinate system $ K $.  At that moment, the clock in $ K $ shows $ t = 15 \, \mu\text{s} $.  At the same time in $ K $, Balázs drives another nail into point $ P_2(3000\,\text{m}, 500\,\text{m}, 150\,\text{m}) $ of the same coordinate system.  (The two events occur simultaneously in $ K $.)  Calculate the time interval between the two events in a reference frame $ K' $ moving at a uniform velocity $ v = 2 \cdot 10^8 \, \text{m/s} $, parallel to the x-axis.

\end{enumerate}

\item Relativistic Dynamics

\begin{enumerate}[label=\arabic*.] 

\item A particle moves with a velocity such that $ \frac{v}{c} = 0.99 $.  Determine the value of $ \frac{m}{m_0} $!

\item Determine the relativistic mass and the velocity of an electron whose kinetic energy is 100\,keV ($ 1.6 \cdot 10^{-14} \, \text{J} $).  The rest mass of the electron is $ 9.11 \cdot 10^{-31} \, \text{kg} $.

\item Calculate the rest energy of an electron, corresponding to its rest mass of $ 9.11 \cdot 10^{-31} \, \text{kg} $, in joules and in electronvolts!  ($ 1 \, \text{MeV} = 1.6 \cdot 10^{-13} \, \text{J} $)

\end{enumerate}

\end{enumerate}

\bigskip

\rightline{Dr.~Gabor~FACSKO}
\rightline{\textit{facsko.gabor@uni-obuda.hu}}

\vfill

\end{document}
